\subsection{The implemented map and its verification output}\label{sec:implemented}
Fix the economy, boundary realization, network classes, optimization budget, and verification procedure, denoted collectively by $\vartheta$. Let $\mathcal P$ be the admissible feedback policies, $\Xi$ the critic parameter space, and $\Omega$ the sampling streams. The implemented NBO is the numerical map
\begin{equation}\label{eq:implemented-r6}
 \mathcal N_\vartheta:\mathcal P\times\Xi\times\Omega
 \longrightarrow\mathcal V_\vartheta\times\mathcal P_\vartheta\times\mathcal C_\vartheta,
 \qquad(a,\xi,\omega)\longmapsto(v,\widetilde a,C).
\end{equation}
The first component is approximate policy evaluation; the second is feasible improvement; the third records verification of the resulting pair. A record identifies the model, parameters, domain, boundary faces, evaluation error, maximization error, coverage, discretization errors, and unchecked assumptions. An unavailable bound is not zero. Failure to obtain a useful certificate is a possible output of the map. With its sampling stream fixed, the implementation is deterministic. We do not assert that a network learns an unrestricted Bellman operator on a function space.

The evaluation--improvement structure is an instance of approximate policy iteration, including its partial-evaluation variants \citep{scherrer2015}. Neural representation does not introduce a new principle of optimality, and monotonicity belongs to the underlying positive-weight scheme, not to gradient descent. The NBO specification joins the block-specific derivative graph to boundary and utility-domain restrictions and to a common economic error account. Sections~\ref{sec:nonquadratic} and~\ref{sec:ndu-neural} implement the differential and positive-weight versions of this map, respectively, with independent verification of their outputs.
